
\[
  \sum_{i=0}^{r+1}(-1)^i\lambda^i(E)f_*(\xi^{\,n-i})=0 .
\]
Le
résultat en découle aussitôt d'après la relation (1.11.3), et l'hypothèse de récurrence.

\textbf{Proposition 5.3.} \uline{Avec les notations et les hypothèses de} 5.1, \uline{on a pour} \uline{tout} $k\ge 0$ \uline{une décomposition en somme directe de} $K^{\bullet}(S)$-\uline{modules} :
\begin{equation}\notag
  \Fil^kK^{\bullet}(X)=\Fil^kK^{\bullet}(S)+\Fil^{k-1}K^{\bullet}(S).x+\ldots+\Fil^{k-r}K^{\bullet}(S).x^r,
\end{equation}
\uline{la filtration considérée étant la filtration de} $\lambda$-\uline{algèbre augmentée définie} \uline{en} V 3.9.1 \uline{et} 3.10.

On remarque d'abord que $x$ est d'augmentation nulle. Par ailleurs, on a :
\[
  \gamma_t(x)=\lambda_{t/(1-t)}(1-\xi^{-1})
  =\left(1+\frac{t}{1-t}\right)\left(1+\frac{\xi^{-1}t}{1-t}\right)^{-1}
  =1/(1-xt) .
\]
Donc $\gamma^i(x)=x^i$ pour tout $i$. Par suite, $K^{\bullet}(X)$ est un quotient de la $K^{\bullet}(S)$-$\lambda$-algèbre engendrée par un générateur $X$ soumis aux relations $\gamma^i(X)=X^i$ pour tout $i$, et l'idéal par lequel on fait le quotient est engendré par la relation (1.13.1). On peut donc appliquer V 4.12 et V 4.13.ii), et on est ramené à montrer que pour tout $n\ge 0$ :
\[
  x^n\in \Fil^nK^{\bullet}(S)+\Fil^{n-1}K^{\bullet}(S).x+\ldots+\Fil^{n-r}K^{\bullet}(S).x^r .
\]

C'est évident pour $n\le r$; on poursuit par récurrence. La relation (1.13.1) donne :
\[
  x^n=-\sum_{i=0}^{r}\gamma^{r+1-i}(\check E-r-1).x^{n-r-1+i},\qquad \text{avec } n\ge r+1;
\]
appliquant l'hypothèse de récurrence aux $x^{n-r-1+i}$, on en déduit le résultat cherché, la somme étant directe puisque $1,\ldots,x^r$ sont indépendants sur $K^{\bullet}(S)$.

